The results thus far are arguably disappointing.
While high test accuracy can be achieved by the boolean model on a variety of
datasets that include both continuous and discrete features before binarization,
such performance is not achieved for our dataset and the rulesets found are inconsistent.

We conduct several additional experience in an attempt to understand why these problems occur.
In \autoref{sec:tissue_status}, we apply the model to subsets of the data.
In particular, we consider the tumors and metastases separately.
This may reveal that tissue status affects TF behavior and may yield better results.
In \autoref{sec:exclusion} we study the effect of excluding TF on performance, to identify
whether there exist TFs or TF combinations that are important capturing PD-L1 expression inside a ruleset.
If performance does not meaningfully decrease for any exclusions,
then this suggests that there do not exist rules that are meaningfully better than others,
and that optimal rulesets should be interpreted with extreme caution.
In \autoref{sec:gene_expression}, we use the model on raw gene expression data and study its behavior.
This may provide an indication that our TF activity data is not suitable for this kind of analysis,
but that the method could be effective on other kinds of biological data.

\section{Tissue status}\label{sec:tissue_status}
It is worth investigating whether there is a difference in results when considering only tumors or only metastases.
If the model performs substantially better on these subsets than on their union,
this is valuable on its own and provides a valuable new insight.
Precisely, it suggests that good rulesets for tumors contradict good rulesets for metastases.

Furthermore, a metastasis tends to have more genetic mutations than a tumor.
If more mutations cause more complicated and more unique behavior within the cell,
it may be harder to deduce general rules governing PD-L1 expression in metastases.
Again, focusing on one or the other independently could be useful in this case.

A related hypothesis is that metastases have higher PD-L1 expression due to their tendency for more mutations.
To investigate this, the distribution of PD-L1 expression per tissue type is shown in \autoref{fig:tissue_status_density_broad}.
Note that cell lines belonging to Papiloma or Normal are disregarded due to a lack of samples.
The density graphs have a similar shape.
The mean PD-L1 expression per tissue status is also similar.
Alternatively, if the hypothesis is true, we expect to see relatively many metastases among samples with high PD-L1 expression.
The tissue status distribution per PD-L1 expression quantile bin is depicted in \autoref{fig:tissue_status_quantile_broad}.
This histogram does not suggest the balance of tumors and metastases changes significantly with PD-L1 expression.
Thus, the results suggest that the hypothesis is false.

Still, there may be a difference in the behavior of the model between tumors and metastases.
To test this, the greedy method was used on the Broad dataset for tumors and metastases separately.
Both the tumors in Broad and the metastases in Broad were treated as separate datasets,
each undergoing independent discretization, and were subjected to the algorithm independently.
We consider $Q_y = 0.5, 0.6$ and $D=2$.
These values of $Q_y$ are the ones that were identified as most promising in \autoref{sec:boolean-results}.
The choice $D=2$ is practical: the greedy method is easier to use,
and for $D=2$, the greedy method is orders of magnitude faster than for higher $D$.
It is not purely practical;
in \autoref{sec:threshold-results} it was shown that performance barely increased when increasing $D$.
However, these results need not extend to this experiment.
It may be that other values of $Q_y$ and higher values of $D$ yield better results on these subsets.

The results are shown in \autoref{tab:greedy_type_tumor_yq_4_mw_0_C_3_D_2} through
\autoref{tab:greedy_type_metastasis_yq_5_mw_0_C_3_D_2}.
We note, without further analysis,
that first rules found in each of the four variants are all qualitatively different.

In all four cases, there is effectively no value in adding more than one rule.
Recall that after the first rule is found, all covered points are stripped away and the best rule for the remaining points
is found through enumeration.
Therefore, there is simply no effective rule for the leftover points.

Accuracy is slightly higher for tumors.
Therefore, the results of this experiment suggest that
PD-L1 expression is slightly more easily captured in a ruleset for tumors.
However, this says nothing about whether the rulesets are biologically meaningful.

Accuracy for metastases is lower than what is observed for the full dataset.
This suggests that contradicting rulesets between tumors and metastases are not the cause of low performance when using the full dataset.



\begin{figure}[h!]
    \centering
    \includesvg[width=0.8\textwidth]{gen/data/tissue_type_density_broad.svg}
    \caption{Density of PD-L1 expression for tissue types Tumor, Metastasis and Unknown in the Broad dataset.}
    \label{fig:tissue_status_density_broad}
\end{figure}


\begin{figure}[h!]
    \centering
    \includesvg[width=0.8\textwidth]{gen/data/tissue_type_quantile_broad.svg}
    \caption{Distribution of tissue type within various quantile bins of PD-L1 expression in the Broad dataset}
    \label{fig:tissue_status_quantile_broad}
\end{figure}

\clearpage

\vspace*{\fill}
\begin{table}[h]
    \centering
    \bwtable{gen/PDL1ex2/tab_greedy_type_tumor_yq_4_mw_0_C_3_D_2.tex}
    \caption{Results of the greedy method for tumors, $Q_y = 0.5$.}
    \label{tab:greedy_type_tumor_yq_4_mw_0_C_3_D_2}
\end{table}

\begin{table}[h]
    \centering
    \bwtable{gen/PDL1ex2/tab_greedy_type_tumor_yq_5_mw_0_C_3_D_2.tex}
    \caption{Results of the greedy method for tumors, $Q_y = 0.6$.}
    \label{tab:greedy_type_tumor_yq_5_mw_0_C_3_D_2}
\end{table}
\vspace*{\fill}

\clearpage

\vspace*{\fill}
\begin{table}[h]
    \centering
    \bwtable{gen/PDL1ex2/tab_greedy_type_metastasis_yq_4_mw_0_C_3_D_2.tex}
    \caption{Results of the greedy method for metastases, $Q_y = 0.5$.}
    \label{tab:greedy_type_metastasis_yq_4_mw_0_C_3_D_2}
\end{table}

\begin{table}[h]
    \centering
    \bwtable{gen/PDL1ex2/tab_greedy_type_metastasis_yq_5_mw_0_C_3_D_2.tex}
    \caption{Results of the greedy method for metastases, $Q_y = 0.6$.}
    \label{tab:greedy_type_metastasis_yq_5_mw_0_C_3_D_2}
\end{table}
\vspace*{\fill}

\clearpage

\section{Transcription Factor Exclusion}\label{sec:exclusion}
Evidently, learning logical rules from bulk transcriptomics is not easy.
Rather than focussing on the rulesets produced and their consistency,
we will now attempt to identify TF that are ``important''.
That is, TF that are essential to classifying PD-L1 expression ``well'' with rulesets.

To do this, we argue by exclusion: if a TF being excluded from the dataset causes performance to drop substantially,
then this TF may be considered important.

Consider $C=1, D=2$.
If the best possible ruleset is, for instance,
\[
    \text{TF A} \wedge \text{TF B},
\]
then removing any TF other than A or B from the dataset will not remove this solution.
Therefore, removing any TF unequal to A or B will have no effect on the maximum attainable performance.
In other words, the only removals that have any effect on performance are either the removal of A or the removal of B\@.

We use this idea to build a tree.
First, the best ruleset is found.
For each TF in the ruleset, a branch is created in which it is excluded.
Then the best ruleset is found and the process repeats.
We consider
\[
    C=1, \quad D=1,2, \quad \text{and} \quad Q_y=0.5,0.6.
\]
Since it is feasible to solve these cases to optimality with the greedy method, this method is used.

For the results, see \autoref{tab:exclusion}.
The reduction in accuracy upon exclusion of a TF is between $0$ and $2$ percentage points.
We consider this to be small,
suggesting that there are no combinations of TF that are truly important to maintain comparable training accuracy.
Still, these tables may provide some insight to a biologist.

\clearpage

\begin{landscape}
    \vspace*{\fill}
    \begin{center}
        \captionof{table}{TF exclusion table as described in \autoref{sec:exclusion}. \\ Case $Q_y=0.5$ (top) and $Q_y=0.6$ (bottom), $D=1$ (left) and $D=2$ (right).}
        \vspace{1.5em}

        \label{tab:exclusion}
        \begin{minipage}{0.1\linewidth}
            \centering
            \begin{tabular}{| c|c|}\hline
\textbf{Depth} & \multicolumn{1}{|>{\columncolor[HTML]{B8E6B8}}c|}{Baseline: } \\
 & \multicolumn{1}{|>{\columncolor[HTML]{B8E6B8}}c|}{0.69} \\ \hline
1 & \multicolumn{1}{|>{\columncolor[HTML]{C4DCB8}}c|}{-STAT2:} \\
 & \multicolumn{1}{|>{\columncolor[HTML]{C4DCB8}}c|}{0.69} \\ \hline
2 & \multicolumn{1}{|>{\columncolor[HTML]{EABFB8}}c|}{-RELA:} \\
 & \multicolumn{1}{|>{\columncolor[HTML]{EABFB8}}c|}{0.66} \\ \hline
3 & \multicolumn{1}{|>{\columncolor[HTML]{F5B8B8}}c|}{-FOSL1:} \\
 & \multicolumn{1}{|>{\columncolor[HTML]{F5B8B8}}c|}{0.66} \\ \hline
\end{tabular}
        \end{minipage}
        \hfill
        \begin{minipage}{0.85\linewidth}
            \centering
            \begin{tabular}{||c|c|c|c|c|c|c|c|}\hline
\multicolumn{8}{|>{\columncolor[HTML]{B8E6B8}}c|}{Baseline: } \\
\multicolumn{8}{|>{\columncolor[HTML]{B8E6B8}}c|}{0.71} \\ \hline
\multicolumn{4}{|>{\columncolor[HTML]{BAE4B8}}c|}{-FOS:} & \multicolumn{4}{|>{\columncolor[HTML]{CDD5B8}}c|}{-STAT2:} \\
\multicolumn{4}{|>{\columncolor[HTML]{BAE4B8}}c|}{0.71} & \multicolumn{4}{|>{\columncolor[HTML]{CDD5B8}}c|}{0.69} \\ \hline
\multicolumn{2}{|>{\columncolor[HTML]{C0DFB8}}c|}{-RELA:} & \multicolumn{2}{|>{\columncolor[HTML]{CDD5B8}}c|}{-STAT2:} & \multicolumn{2}{|>{\columncolor[HTML]{E5C3B8}}c|}{-RELA:} & \multicolumn{2}{|>{\columncolor[HTML]{CFD3B8}}c|}{-SOX2:} \\
\multicolumn{2}{|>{\columncolor[HTML]{C0DFB8}}c|}{0.70} & \multicolumn{2}{|>{\columncolor[HTML]{CDD5B8}}c|}{0.69} & \multicolumn{2}{|>{\columncolor[HTML]{E5C3B8}}c|}{0.68} & \multicolumn{2}{|>{\columncolor[HTML]{CFD3B8}}c|}{0.69} \\ \hline
\multicolumn{1}{|>{\columncolor[HTML]{C4DCB8}}c|}{-NFKB1:} & \multicolumn{1}{|>{\columncolor[HTML]{E5C3B8}}c|}{-STAT2:} & \multicolumn{1}{|>{\columncolor[HTML]{E5C3B8}}c|}{-RELA:} & \multicolumn{1}{|>{\columncolor[HTML]{D2D2B8}}c|}{-SOX2:} & \multicolumn{1}{|>{\columncolor[HTML]{F5B8B8}}c|}{-FOSL1:} & \multicolumn{1}{|>{\columncolor[HTML]{EAC0B8}}c|}{-NFKB1:} & \multicolumn{1}{|>{\columncolor[HTML]{D2D2B8}}c|}{-FOS:} & \multicolumn{1}{|>{\columncolor[HTML]{E5C3B8}}c|}{-RELA:} \\
\multicolumn{1}{|>{\columncolor[HTML]{C4DCB8}}c|}{0.70} & \multicolumn{1}{|>{\columncolor[HTML]{E5C3B8}}c|}{0.68} & \multicolumn{1}{|>{\columncolor[HTML]{E5C3B8}}c|}{0.68} & \multicolumn{1}{|>{\columncolor[HTML]{D2D2B8}}c|}{0.69} & \multicolumn{1}{|>{\columncolor[HTML]{F5B8B8}}c|}{0.67} & \multicolumn{1}{|>{\columncolor[HTML]{EAC0B8}}c|}{0.68} & \multicolumn{1}{|>{\columncolor[HTML]{D2D2B8}}c|}{0.69} & \multicolumn{1}{|>{\columncolor[HTML]{E5C3B8}}c|}{0.68} \\ \hline
\end{tabular}
        \end{minipage}

        \vspace{1.5em}

        \begin{minipage}{0.1\linewidth}
            \centering
            \begin{tabular}{| c|c|}\hline
\textbf{Depth} & \multicolumn{1}{|>{\columncolor[HTML]{B8E6B8}}c|}{Baseline: } \\
 & \multicolumn{1}{|>{\columncolor[HTML]{B8E6B8}}c|}{0.70} \\ \hline
1 & \multicolumn{1}{|>{\columncolor[HTML]{C1DEB8}}c|}{-FOSL1:} \\
 & \multicolumn{1}{|>{\columncolor[HTML]{C1DEB8}}c|}{0.70} \\ \hline
2 & \multicolumn{1}{|>{\columncolor[HTML]{D8CDB8}}c|}{-STAT2:} \\
 & \multicolumn{1}{|>{\columncolor[HTML]{D8CDB8}}c|}{0.69} \\ \hline
3 & \multicolumn{1}{|>{\columncolor[HTML]{F5B8B8}}c|}{-RELA:} \\
 & \multicolumn{1}{|>{\columncolor[HTML]{F5B8B8}}c|}{0.69} \\ \hline
\end{tabular}
        \end{minipage}
        \hfill
        \begin{minipage}{0.85\linewidth}
            \centering
            \begin{tabular}{||c|c|c|c|c|c|c|c|}\hline
\multicolumn{8}{|>{\columncolor[HTML]{B8E6B8}}c|}{Baseline: } \\
\multicolumn{8}{|>{\columncolor[HTML]{B8E6B8}}c|}{0.74} \\ \hline
\multicolumn{4}{|>{\columncolor[HTML]{CAD8B8}}c|}{-FOSL1:} & \multicolumn{4}{|>{\columncolor[HTML]{C8D9B8}}c|}{-STAT2:} \\
\multicolumn{4}{|>{\columncolor[HTML]{CAD8B8}}c|}{0.73} & \multicolumn{4}{|>{\columncolor[HTML]{C8D9B8}}c|}{0.73} \\ \hline
\multicolumn{2}{|>{\columncolor[HTML]{D6CFB8}}c|}{-RELA:} & \multicolumn{2}{|>{\columncolor[HTML]{ECBEB8}}c|}{-STAT2:} & \multicolumn{2}{|>{\columncolor[HTML]{ECBEB8}}c|}{-FOSL1:} & \multicolumn{2}{|>{\columncolor[HTML]{CAD8B8}}c|}{-RELA:} \\
\multicolumn{2}{|>{\columncolor[HTML]{D6CFB8}}c|}{0.72} & \multicolumn{2}{|>{\columncolor[HTML]{ECBEB8}}c|}{0.71} & \multicolumn{2}{|>{\columncolor[HTML]{ECBEB8}}c|}{0.71} & \multicolumn{2}{|>{\columncolor[HTML]{CAD8B8}}c|}{0.73} \\ \hline
\multicolumn{1}{|>{\columncolor[HTML]{D6CFB8}}c|}{-FOS:} & \multicolumn{1}{|>{\columncolor[HTML]{F5B8B8}}c|}{-STAT2:} & \multicolumn{1}{|>{\columncolor[HTML]{F5B8B8}}c|}{-RELA:} & \multicolumn{1}{|>{\columncolor[HTML]{ECBEB8}}c|}{-SOX2:} & \multicolumn{1}{|>{\columncolor[HTML]{F5B8B8}}c|}{-RELA:} & \multicolumn{1}{|>{\columncolor[HTML]{ECBEB8}}c|}{-SOX2:} & \multicolumn{1}{|>{\columncolor[HTML]{D6CFB8}}c|}{-FOS:} & \multicolumn{1}{|>{\columncolor[HTML]{F5B8B8}}c|}{-FOSL1:} \\
\multicolumn{1}{|>{\columncolor[HTML]{D6CFB8}}c|}{0.72} & \multicolumn{1}{|>{\columncolor[HTML]{F5B8B8}}c|}{0.70} & \multicolumn{1}{|>{\columncolor[HTML]{F5B8B8}}c|}{0.70} & \multicolumn{1}{|>{\columncolor[HTML]{ECBEB8}}c|}{0.71} & \multicolumn{1}{|>{\columncolor[HTML]{F5B8B8}}c|}{0.70} & \multicolumn{1}{|>{\columncolor[HTML]{ECBEB8}}c|}{0.71} & \multicolumn{1}{|>{\columncolor[HTML]{D6CFB8}}c|}{0.72} & \multicolumn{1}{|>{\columncolor[HTML]{F5B8B8}}c|}{0.70} \\ \hline
\end{tabular}
        \end{minipage}
    \end{center}
    \vspace*{\fill}
\end{landscape}

\section{Gene expression as input}\label{sec:gene_expression}
The use of the DecoupleR model (\cite{badia-i-mompel_decoupler_2022, badia-i-mompel_transcription_nodate}) raises some suspicion.
The model itself assumes that TF act individually to regulate a set of target genes.
A prior knowledge matrix, in our case based on CollecTRI (\cite{muller-dott_expanding_2023}),
describes which TF are known to up- or downregulate which genes.
The gene expressions of all the TF's target genes are then use to ``reverse engineer'' (estimate) the TF activity.
DecoupleR's output is taken as an undisputed truth onwards.
It is unclear whether combinatoric effects can be observed from output of a model that does not assume combinatoric behavior.

Furthermore, the rules we have found may fail to capture real-life phenomena
and instead reflect information from the prior knowledge matrix.
Due to the construction of the DecoupleR model, TFs that target similar genes will have similar activity estimates.
This may partly explain the high cross-correlations observed in \autoref{fig:corr_activity}.
As a result, even if TF activity differs in reality,
the DecoupleR model can make TF activities hard to distinguish.
This, in turn, can make finding meaningful rules difficult.
Instead, the rules may represent the prior knowledge matrix, meaning nothing of interest is learned.

As a curiosity, we use the threshold model (\autoref{ch:threshold-model}) on the raw gene expression data instead.
We do not use raw read counts but normalized TPM counts, meant to be comparable across samples.
While this is not expected to yield a result with a realistic interpretation,
it can possibly give insight in how the model behaves.

\autoref{fig:tpm_accuracy_mw_0.svg} shows the results, in the same way as \autoref{fig:threshold_accuracy_mw_0.svg} does.
The interpretable results for $C=1$ and $D=1,2,3$ are given in \autoref{tab:tpm_interpretable_results_C_1}.
The interpretable results for $C=2$ are given in \autoref{tab:tpm_interpretable_results_C_2}.

It follows that the accuracy achieved is higher compared to the experiments with TF activity estimates.
Furthermore, qualitatively different results are observed.
This indicates that gene expression is indeed not equivalent to TF activity.

\begin{figure}[h]
    \centering
    \includesvg[width=0.7\textwidth]{gen/PDL1ex3/PDL1ex3_fig_accuracy_mw_0.svg}
    \caption{Accuracy of results for the full Broad dataset, using gene expression as input, for $C=1,2,3,4,5$, $D=1,2,3,4$ and $Q_y=0.5,0.6,0.7,0.8,0.9.$}
    \label{fig:tpm_accuracy_mw_0.svg}
\end{figure}

\begin{table}[h]
    \centering
    \bwtable{gen/PDL1ex3/PDL1ex3_tab_interpretable_mw_0_C_1}
    \caption{Results for the full Broad dataset, using gene expression as input, for $C=1$ and small rules ($D=1,2,3$).}
    \label{tab:tpm_interpretable_results_C_1}
\end{table}

\begin{table}[h]
    \bwtable{gen/PDL1ex3/PDL1ex3_tab_interpretable_mw_0_C_2}
    \caption{Results for the full Broad dataset, using gene expression input, for $C=2$ and small rules ($D=1,2,3$).}
    \label{tab:tpm_interpretable_results_C_2}
\end{table}

Train accuracy is not the full story.
It remains unclear whether the results are biologically meaningful.
To study whether these results are meaningful in a mathematical sense,
we perform the exclusion experiment on this dataset too.
If there is a substantial drop in performance upon exclusion of a TF or a combination of TFs,
then this suggests that there do exist rules for the gene expression dataset that are more important than others.

The TF exclusion tables for $Q_y=0.5,0.6$ and $D=2$, the same parameters used earlier,
can be found in \autoref{tab:exclusion_tpm}.
Accuracy decreases remain small (between $0$ and $2$ percentage points), with two notable exceptions:
FOSL1 and IRF1.
In particular, for $Q_y=0.6$, excluding FOSL1 and IRF1 results in a substantial performance drop.
In the case $D=1$, exclusion of FOSL1 causes maximum attainable train accuracy to drop $9$ percentage points.
Then removing IRF1 has little effect.
More notably, for $D=2$, excluding FOSL1 and IRF1 causes an $8$ percentage point drop in the left branch and a
$9$ percentage point drop in the right branch, which we consider a substantial decrease.

This suggests that FOSL1 is an important TF with respect to this dataset, and that FOSL1 and IRF1 are an important
combination.
Furthermore, it shows that an exclusion experiment can yield new insights into TF combinations.
By extension, it strengthens the idea that TF activity estimates, as used in previous chapters,
are not suitable for inferring rulesets.

\clearpage
\begin{landscape}
    \vspace*{\fill}
    \begin{center}
        \captionof{table}{TF exclusion table as described in \autoref{sec:exclusion} for raw gene expression data,
            as described in \autoref{sec:gene_expression}.
        \\ Case $Q_y=0.5$ (top) and $Q_y=0.6$ (bottom), $D=1$ (left) and $D=2$ (right).}
        \vspace{1.5em}

        \label{tab:exclusion_tpm}
        \begin{minipage}{0.1\linewidth}
            \centering
            \begin{tabular}{||c|}\hline
\multicolumn{1}{|>{\columncolor[HTML]{B8E6B8}}c|}{Baseline: } \\
\multicolumn{1}{|>{\columncolor[HTML]{B8E6B8}}c|}{0.71} \\ \hline
\multicolumn{1}{|>{\columncolor[HTML]{F0BBB8}}c|}{-FOSL1:} \\
\multicolumn{1}{|>{\columncolor[HTML]{F0BBB8}}c|}{0.65} \\ \hline
\multicolumn{1}{|>{\columncolor[HTML]{F1BAB8}}c|}{-IRF1:} \\
\multicolumn{1}{|>{\columncolor[HTML]{F1BAB8}}c|}{0.65} \\ \hline
\multicolumn{1}{|>{\columncolor[HTML]{F5B8B8}}c|}{-STAT1:} \\
\multicolumn{1}{|>{\columncolor[HTML]{F5B8B8}}c|}{0.64} \\ \hline
\end{tabular}
        \end{minipage}
        \hfill
        \begin{minipage}{0.85\linewidth}
            \centering
            \begin{tabular}{||c|c|c|c|c|c|c|c|}\hline
\multicolumn{8}{|>{\columncolor[HTML]{B8E6B8}}c|}{Baseline: } \\
\multicolumn{8}{|>{\columncolor[HTML]{B8E6B8}}c|}{0.73} \\ \hline
\multicolumn{4}{|>{\columncolor[HTML]{DACCB8}}c|}{-FOSL1:} & \multicolumn{4}{|>{\columncolor[HTML]{B9E5B8}}c|}{-STAT4:} \\
\multicolumn{4}{|>{\columncolor[HTML]{DACCB8}}c|}{0.70} & \multicolumn{4}{|>{\columncolor[HTML]{B9E5B8}}c|}{0.73} \\ \hline
\multicolumn{2}{|>{\columncolor[HTML]{E7C2B8}}c|}{-IRF1:} & \multicolumn{2}{|>{\columncolor[HTML]{DEC8B8}}c|}{-NFATC1:} & \multicolumn{2}{|>{\columncolor[HTML]{DACCB8}}c|}{-FOSL1:} & \multicolumn{2}{|>{\columncolor[HTML]{BCE2B8}}c|}{-NFATC1:} \\
\multicolumn{2}{|>{\columncolor[HTML]{E7C2B8}}c|}{0.69} & \multicolumn{2}{|>{\columncolor[HTML]{DEC8B8}}c|}{0.70} & \multicolumn{2}{|>{\columncolor[HTML]{DACCB8}}c|}{0.70} & \multicolumn{2}{|>{\columncolor[HTML]{BCE2B8}}c|}{0.73} \\ \hline
\multicolumn{1}{|>{\columncolor[HTML]{F5B8B8}}c|}{-STAT1:} & \multicolumn{1}{|>{\columncolor[HTML]{F0BBB8}}c|}{-NFATC1:} & \multicolumn{1}{|>{\columncolor[HTML]{F0BBB8}}c|}{-IRF1:} & \multicolumn{1}{|>{\columncolor[HTML]{E0C7B8}}c|}{-SPI1:} & \multicolumn{1}{|>{\columncolor[HTML]{E7C2B8}}c|}{-IRF1:} & \multicolumn{1}{|>{\columncolor[HTML]{DEC8B8}}c|}{-NFATC1:} & \multicolumn{1}{|>{\columncolor[HTML]{DEC8B8}}c|}{-FOSL1:} & \multicolumn{1}{|>{\columncolor[HTML]{C3DDB8}}c|}{-IRF1:} \\
\multicolumn{1}{|>{\columncolor[HTML]{F5B8B8}}c|}{0.68} & \multicolumn{1}{|>{\columncolor[HTML]{F0BBB8}}c|}{0.68} & \multicolumn{1}{|>{\columncolor[HTML]{F0BBB8}}c|}{0.68} & \multicolumn{1}{|>{\columncolor[HTML]{E0C7B8}}c|}{0.70} & \multicolumn{1}{|>{\columncolor[HTML]{E7C2B8}}c|}{0.69} & \multicolumn{1}{|>{\columncolor[HTML]{DEC8B8}}c|}{0.70} & \multicolumn{1}{|>{\columncolor[HTML]{DEC8B8}}c|}{0.70} & \multicolumn{1}{|>{\columncolor[HTML]{C3DDB8}}c|}{0.72} \\ \hline
\end{tabular}
        \end{minipage}

        \vspace{1.5em}

        \begin{minipage}{0.1\linewidth}
            \centering
            \begin{tabular}{||c|}\hline
\multicolumn{1}{|>{\columncolor[HTML]{B8E6B8}}c|}{Baseline: } \\
\multicolumn{1}{|>{\columncolor[HTML]{B8E6B8}}c|}{0.74} \\ \hline
\multicolumn{1}{|>{\columncolor[HTML]{E9C0B8}}c|}{-FOSL1:} \\
\multicolumn{1}{|>{\columncolor[HTML]{E9C0B8}}c|}{0.65} \\ \hline
\multicolumn{1}{|>{\columncolor[HTML]{F0BBB8}}c|}{-IRF1:} \\
\multicolumn{1}{|>{\columncolor[HTML]{F0BBB8}}c|}{0.64} \\ \hline
\multicolumn{1}{|>{\columncolor[HTML]{F5B8B8}}c|}{-STAT4:} \\
\multicolumn{1}{|>{\columncolor[HTML]{F5B8B8}}c|}{0.63} \\ \hline
\end{tabular}
        \end{minipage}
        \hfill
        \begin{minipage}{0.85\linewidth}
            \centering
            \begin{tabular}{||c|c|c|c|c|c|c|c|}\hline
\multicolumn{8}{|>{\columncolor[HTML]{B8E6B8}}c|}{Baseline: } \\
\multicolumn{8}{|>{\columncolor[HTML]{B8E6B8}}c|}{0.76} \\ \hline
\multicolumn{4}{|>{\columncolor[HTML]{D7CEB8}}c|}{-FOSL1:} & \multicolumn{4}{|>{\columncolor[HTML]{BEE1B8}}c|}{-NFATC1:} \\
\multicolumn{4}{|>{\columncolor[HTML]{D7CEB8}}c|}{0.71} & \multicolumn{4}{|>{\columncolor[HTML]{BEE1B8}}c|}{0.75} \\ \hline
\multicolumn{2}{|>{\columncolor[HTML]{E7C1B8}}c|}{-IRF1:} & \multicolumn{2}{|>{\columncolor[HTML]{D8CDB8}}c|}{-NFATC1:} & \multicolumn{2}{|>{\columncolor[HTML]{D8CDB8}}c|}{-FOSL1:} & \multicolumn{2}{|>{\columncolor[HTML]{BFE0B8}}c|}{-IRF1:} \\
\multicolumn{2}{|>{\columncolor[HTML]{E7C1B8}}c|}{0.68} & \multicolumn{2}{|>{\columncolor[HTML]{D8CDB8}}c|}{0.71} & \multicolumn{2}{|>{\columncolor[HTML]{D8CDB8}}c|}{0.71} & \multicolumn{2}{|>{\columncolor[HTML]{BFE0B8}}c|}{0.75} \\ \hline
\multicolumn{1}{|>{\columncolor[HTML]{E9C0B8}}c|}{-STAT1:} & \multicolumn{1}{|>{\columncolor[HTML]{F5B8B8}}c|}{-NFATC1:} & \multicolumn{1}{|>{\columncolor[HTML]{F5B8B8}}c|}{-IRF1:} & \multicolumn{1}{|>{\columncolor[HTML]{DCCAB8}}c|}{-MYC:} & \multicolumn{1}{|>{\columncolor[HTML]{F5B8B8}}c|}{-IRF1:} & \multicolumn{1}{|>{\columncolor[HTML]{DCCAB8}}c|}{-MYC:} & \multicolumn{1}{|>{\columncolor[HTML]{F5B8B8}}c|}{-FOSL1:} & \multicolumn{1}{|>{\columncolor[HTML]{C3DDB8}}c|}{-STAT4:} \\
\multicolumn{1}{|>{\columncolor[HTML]{E9C0B8}}c|}{0.68} & \multicolumn{1}{|>{\columncolor[HTML]{F5B8B8}}c|}{0.66} & \multicolumn{1}{|>{\columncolor[HTML]{F5B8B8}}c|}{0.66} & \multicolumn{1}{|>{\columncolor[HTML]{DCCAB8}}c|}{0.70} & \multicolumn{1}{|>{\columncolor[HTML]{F5B8B8}}c|}{0.66} & \multicolumn{1}{|>{\columncolor[HTML]{DCCAB8}}c|}{0.70} & \multicolumn{1}{|>{\columncolor[HTML]{F5B8B8}}c|}{0.66} & \multicolumn{1}{|>{\columncolor[HTML]{C3DDB8}}c|}{0.74} \\ \hline
\end{tabular}
        \end{minipage}
    \end{center}
    \vspace*{\fill}
\end{landscape}


